\documentclass[../main.tex]{subfiles}

\graphicspath{{../image/}{../figure/}}

\begin{document}

\chapter{熵与经典解的稳定性}\label{chapter:ch05-熵与经典解稳定性}

连续介质物理的一条基本信条是：热力学第二定律本质上是一条稳定性陈述。在前述各章讨论的例子中，第二定律表现为伴随平衡律的存在——对经典解以等式恒等满足，而对平衡律方程组的弱解则作为热力学可容许性不等式约束施加。在本书关于双曲平衡律方程组理论的阐述中，一个反复出现的主题是：伴随平衡律以各种形式诱导出稳定性。在此，读者将一瞥熵不等式对经典解稳定性的蕴含。

我们将证明：当平衡律方程组被一个由凸熵诱导的伴随平衡律装备时，初值问题在经典解的框架下局部适定：充分光滑的初值在一个极大时间区间上生成经典解，该区间通常为有限时长。然而，在由松弛或其他耗散机制诱导的阻尼存在时，且当初值充分小时，经典解在时间上整体存在。经典解唯一，并连续依赖于其初值——不仅在经典解类中如此，甚至在更广泛的、将伴随平衡律作为不等式可容许性约束满足的弱解类中亦然。

即使熵不凸，在以下两种情形中也可建立类似的存在性与稳定性结果：(a) 熵仅在状态空间某一锥的方向上凸，但方程组装备有特殊的伴随平衡律，称为对合（involution），其存在补偿了互补方向上凸性的缺失；或 (b) 方程组装备有互补熵且主熵为多凸（polyconvex）。这种结构出现于弹性动力学与电磁学中。

本章末尾将简要讨论初边值问题经典解的存在性。

从分析技巧的角度看，本章呈现了双曲平衡律方程组理论中可由线性理论方法处理的那些方面，即能量估计与 Fourier 分析。

\subfile{Chapter-05/section-01.tex}

\subfile{Chapter-05/section-02.tex}

\subfile{Chapter-05/section-03.tex}

\subfile{Chapter-05/section-04.tex}

\subfile{Chapter-05/section-05.tex}

\subfile{Chapter-05/section-06.tex}

\subfile{Chapter-05/section-07.tex}

\end{document}
